\chapter{Praktikum 2: Sistem Persamaan Linear Iteratif}

\CP{Mahasiswa mampu membangun algoritma Jacobi dan Gauss-Seidel dari bentuk persamaan skalar, kemudian menggunakan notasi matriks, mengevaluasi residual, galat, dan konvergensi.}

\Target{Implementasi Jacobi dan Gauss-Seidel, tabel iterasi, solusi pembanding, residual, dan grafik \texttt{errs}.}

\section{Mulai dari persamaan, bukan matriks}

Pertimbangkan tiga persamaan

\begin{align}
a_{11}x_1+a_{12}x_2+a_{13}x_3 &= b_1,\\
a_{21}x_1+a_{22}x_2+a_{23}x_3 &= b_2,\\
a_{31}x_1+a_{32}x_2+a_{33}x_3 &= b_3.
\end{align}

Masing-masing persamaan diselesaikan terhadap variabel diagonalnya:

\begin{align}
x_1 &=
\frac{b_1-a_{12}x_2-a_{13}x_3}{a_{11}},\\
x_2 &=
\frac{b_2-a_{21}x_1-a_{23}x_3}{a_{22}},\\
x_3 &=
\frac{b_3-a_{31}x_1-a_{32}x_2}{a_{33}}.
\end{align}

Dari bentuk inilah Jacobi dan Gauss-Seidel dibangun.

\section{Formula Jacobi}

Jacobi menggunakan seluruh nilai dari iterasi lama:

\begin{align}
x_1^{(k+1)}
&=
\frac{b_1-a_{12}x_2^{(k)}-a_{13}x_3^{(k)}}{a_{11}},\\
x_2^{(k+1)}
&=
\frac{b_2-a_{21}x_1^{(k)}-a_{23}x_3^{(k)}}{a_{22}},\\
x_3^{(k+1)}
&=
\frac{b_3-a_{31}x_1^{(k)}-a_{32}x_2^{(k)}}{a_{33}}.
\end{align}

Untuk $n$ variabel,

\begin{equation}
x_i^{(k+1)}
=
\frac{1}{a_{ii}}
\left(
b_i-\sum_{\substack{j=1\\j\ne i}}^n a_{ij}x_j^{(k)}
\right).
\end{equation}

\section{Formula Gauss-Seidel}

Gauss-Seidel langsung memakai nilai baru yang sudah tersedia:

\begin{align}
x_1^{(k+1)}
&=
\frac{b_1-a_{12}x_2^{(k)}-a_{13}x_3^{(k)}}{a_{11}},\\
x_2^{(k+1)}
&=
\frac{b_2-a_{21}x_1^{(k+1)}-a_{23}x_3^{(k)}}{a_{22}},\\
x_3^{(k+1)}
&=
\frac{b_3-a_{31}x_1^{(k+1)}-a_{32}x_2^{(k+1)}}{a_{33}}.
\end{align}

Untuk $n$ variabel,

\begin{equation}
x_i^{(k+1)}
=
\frac{1}{a_{ii}}
\left[
b_i
-\sum_{j=1}^{i-1}a_{ij}x_j^{(k+1)}
-\sum_{j=i+1}^{n}a_{ij}x_j^{(k)}
\right].
\end{equation}

Perbedaan utama:

\begin{itemize}
\item Jacobi: semua komponen memakai iterasi $k$.
\item Gauss-Seidel: komponen yang sudah dihitung langsung memakai iterasi $k+1$.
\end{itemize}

\section{Baru kemudian notasi matriks}

Sistem dapat diringkas menjadi

\begin{equation}
A\mathbf{x}=\mathbf{b}.
\end{equation}

Notasi matriks berguna untuk implementasi umum, tetapi algoritma di atas tetap berasal dari persamaan skalar.

\section{Galat iterasi dan residual}

Gunakan

\begin{equation}
e_a^{(k)}
=
\max_i
\frac{|x_i^{(k+1)}-x_i^{(k)}|}
{\max(1,|x_i^{(k+1)}|)}.
\end{equation}

Residual adalah

\begin{equation}
\mathbf{r}^{(k)}
=
\mathbf{b}-A\mathbf{x}^{(k)},
\end{equation}

dan salah satu ukuran residual adalah

\begin{equation}
r_k=\|\mathbf{r}^{(k)}\|_\infty.
\end{equation}

\section{Contoh}

Gunakan

\begin{align}
10x_1-x_2+2x_3&=6,\\
-x_1+11x_2-x_3&=25,\\
2x_1-x_2+10x_3&=-11.
\end{align}

Cari solusi eksak/analitik terlebih dahulu sebagai pembanding, lalu lakukan Jacobi dan Gauss-Seidel dari tebakan $(0,0,0)^T$.

\section{Program MATLAB}

\texttt{spl\_demo.m}

\lstinputlisting[language=Matlab]{matlab/spl_demo.m}

\section{Grafik konvergensi}

Simpan galat setiap iterasi pada vektor \texttt{errs}.

\begin{lstlisting}[language=Matlab]
semilogy(1:numel(errs),errs,'o-')
grid on
xlabel('Iterasi')
ylabel('e_a')
title('Grafik konvergensi')
\end{lstlisting}

Jika membandingkan Jacobi dan Gauss-Seidel, gambarkan kedua kurva pada satu grafik.

\section{Latihan}

\begin{enumerate}
\item Turunkan formula iterasi dari sistem persamaan, jangan mulai dari notasi matriks.
\item Tentukan solusi eksak/analitik sebagai pembanding.
\item Lakukan Jacobi dan Gauss-Seidel.
\item Hitung galat iterasi dan residual.
\item Simpan sejarah galat pada \texttt{errs}.
\item Buat grafik konvergensi kedua metode.
\item Jelaskan metode mana yang lebih cepat pada kasus tersebut.
\end{enumerate}

\section{Latihan terintegrasi dari Lembar Kerja}

\begin{enumerate}
\item[\textbf{B3.1 (L1)}] Tuliskan sistem \(10x_1-x_2+2x_3=6\), \(-x_1+11x_2-x_3=25\), \(2x_1-x_2+10x_3=-11\) ke bentuk \(A\mathbf{x}=\mathbf{b}\).
\item[\textbf{B3.2 (L2)}] Tentukan solusi eksak/analitik SPL tersebut.
\item[\textbf{B3.3 (L3)}] Gunakan Jacobi dari \((0,0,0)^T\), lakukan 3 iterasi, dan hitung galat setiap komponen.
\item[\textbf{B3.4 (L4)}] Gunakan Gauss-Seidel dari tebakan yang sama, lakukan 3 iterasi, lalu bandingkan dengan solusi eksak.
\item[\textbf{B3.5 (L5)}] Untuk model tiga reaktor, jelaskan pembentukan SPL dari neraca massa, susun bentuk matriksnya, dan tentukan solusi eksaknya.
\end{enumerate}
